% =====================================================================
\section{Introduction}
% =====================================================================

\subsection{Background}

The project serves parking facilities: car parks of apartment and office buildings, shopping-center
garages and paid lots for cars and motorbikes. Staff record each vehicle at the entry and exit lanes,
look a vehicle up when a customer loses a ticket or disputes a fee, and reconcile the records at the
end of a shift. Automatic License Plate Recognition (ALPR) turns camera images into plate text and
performs well when the plate is large and sharp \cite{laroca2018robust}. Many facilities, however,
rely on low-cost overview cameras placed some distance from the lane and keep compressed recordings,
so the plate may cover only a few dozen pixels. In those cases staff fall back to reading the footage
themselves.

The ICPR 2026 Low-Resolution License Plate Recognition (LRLPR) benchmark \cite{laroca2026icpr}
captures this setting. It provides 20{,}000 training \emph{tracks}. Each track holds five
low-resolution (LR) crops of the same plate, about $46\times19$ pixels on average, five
high-resolution (HR) crops and the plate text. Plates follow two Brazilian layouts, \code{LLLNNNN}
and \code{LLLNLNN}, and a track counts as correct only when all seven characters match. With five
to seven pixels of width per character, a single frame often leaves some glyphs ambiguous.

The author's model MF-SR-OCR \cite{nguyen2026mfsrocr} restores and reads these tracks by combining
the five frames. Today it exists as research code: training scripts, evaluation scripts and
checkpoints. This capstone turns it into a web application that parking staff can use. The benchmark
plates are Brazilian and the model is trained on them, so the system is built and evaluated on that
domain. Vietnamese plates follow other layouts, often on two rows, and supporting them needs labeled
Vietnamese data; this is outside the present scope (Section~1.4).

\textbf{Stakeholders.} Table~\ref{tab:stakeholders} lists who is affected and what each group needs.
Fig.~\ref{fig:usecase} shows the use cases in scope.

\begin{table}[tbp]
\centering
\caption{Stakeholders and their needs.}
\label{tab:stakeholders}
\begin{tabularx}{\textwidth}{@{}P{30mm} L@{}}
\toprule
Stakeholder & Need \\
\midrule
Parking attendants and operators (primary users) & Read plates at the entry and exit lanes or from recorded footage; know which readings to trust; correct the rest with as little typing as possible. \\
Parking managers & Process the footage of a shift or a day as one batch, export the records, and settle lost-ticket or fee disputes from them. \\
Vehicle owners (indirect) & Their vehicle is identified correctly at exit, and a dispute can be checked against a stored record. \\
System administrators & Install, update and back up the service on CPU-only machines. \\
ML engineer (the author) & Swap a model or threshold without code changes; trace every result to a model version. \\
Dataset owners (UFPR) & Images are used only for non-commercial academic work and are not redistributed \cite{icpr26license}. \\
\bottomrule
\end{tabularx}
\end{table}

\fig{usecase.pdf}{Use cases in scope. Operators recognize, review and export; the administrator
runs the stack; the ML engineer manages the model and the threshold $\tau$.}{fig:usecase}

\subsection{Practical Problem Statement and User Pain-Point}

Parking staff currently have no tool that uses all the frames in which a vehicle appears, indicates
which automatic readings are reliable, and helps them correct the others. Fig.~\ref{fig:workflow}
compares the current workflow with the proposed one. The problem breaks down into five concrete
pain points.

\begin{description}
  \item[P1. Accuracy.] On the 999 validation tracks, a single-frame CRNN reads $74.45\%$ of tracks
        correctly and the multi-frame CRNN+STN baseline, retrained on this codebase, reads
        $75.78\%$. About one plate in four is wrong.
  \item[P2. No reliability signal.] Results come without a usable confidence, so the operator has
        to check every plate to find the wrong ones.
  \item[P3. Slow correction.] A wrong plate is re-typed from a blurry crop, with no hint of which
        character is doubtful and no list of likely alternatives.
  \item[P4. Not deployable.] The model runs only as Python scripts. There is no API, no batch mode,
        no record of corrections, and the target parking sites run ordinary PCs without a GPU.
  \item[P5. Concurrent load.] Batch jobs write results continuously while reviewers read the queue.
        A file-locking database such as SQLite allows one writer at a time \cite{sqlite_lock}, so
        reviewers would wait behind the batch.
\end{description}

\fig{workflow.pdf}{Current workflow (a) and proposed workflow with confidence-based routing (b).}{fig:workflow}

\subsection{Research Objectives and Questions}

Table~\ref{tab:objectives} states the engineering objectives in SMART form. Week~1 starts on
October 5, 2026.

\begin{table}[tbp]
\centering
\caption{SMART objectives.}
\label{tab:objectives}
\begin{tabularx}{\textwidth}{@{}>{\bfseries}l P{27mm} L l@{}}
\toprule
ID & Objective & Measurable target & Due \\
\midrule
O1 & Validate the model & Mean $\pm$ std over 3 seeds $\ge 79\%$ on validation; paired McNemar test and bootstrap 95\% CI against the baseline; one run on a sealed 1{,}000-track test set & W3, W6 \\
O2 & Selective prediction & Threshold $\tau$ that gives $\ge 95\%$ precision on auto-accepted validation tracks; coverage reported & W5 \\
O3 & REST API & 0 divergent tracks of 999 between API and offline model; p95 latency $\le 1$\,s per track on CPU & W7 \\
O4 & Web interface & Single-track, batch and review screens usable on desktop and mobile browsers & W8--9 \\
O5 & Reliable operation & 1{,}000-track batch with no error; 30/30 end-to-end cases pass; RPO $\le 6$\,h, RTO $< 5$\,min & W9--12 \\
O6 & User value & 5-user study with SUS $\ge 68$ \cite{brooke1996sus}; $\ge 30\%$ less review time per flagged track than manual reading & W10 \\
\bottomrule
\end{tabularx}
\end{table}

Three research questions guide the evaluation.

\begin{description}
  \item[RQ1.] Does the multi-frame model improve on the CRNN+STN baseline on data that played no
        part in model selection?
  \item[RQ2.] What share of tracks can be accepted automatically while keeping precision at or
        above $95\%$?
  \item[RQ3.] Do the review aids (super-resolved frames, per-character confidence and top-5
        suggestions) shorten the time a reviewer needs per flagged track?
\end{description}

\subsection{Scope}

Table~\ref{tab:scope} separates what will be built and tested from what is left out. The model is
trained on Brazilian plates, so the system is described as working on that domain only.

\begin{table}[tbp]
\centering
\caption{Project scope.}
\label{tab:scope}
\begin{tabularx}{\textwidth}{@{}L L@{}}
\toprule
\textbf{In-scope components} & \textbf{Out-of-scope components} \\
\midrule
MF-SR-OCR inference with greedy and layout-constrained decoding &
Plate detection and tracking in raw video (stretch goal only); gate hardware such as barriers, ticket machines or RFID readers \\
Selective prediction: track confidence, threshold $\tau$, review queue &
Non-Brazilian plates (10 Vietnamese plates tested qualitatively to document the limitation) \\
FastAPI REST API: single track, ZIP batch, review, suggestions, CSV export &
In-app user accounts (identity is delegated to Cloudflare Access) \\
PostgreSQL storage of results and review history (text only, no images) &
GPU or edge deployment, INT8 quantization, Kubernetes \\
React web app with single-track, batch and review screens &
Live camera streams and integration with police or VMS systems \\
Docker Compose deployment on one CPU host, Cloudflare Tunnel, backups &
Storing or publishing dataset images \\
Automated tests, CI/CD pipeline, API and user documentation &
New model research beyond a minimal ablation \\
Roles: operator/reviewer, administrator, ML engineer & \\
\bottomrule
\end{tabularx}
\end{table}
